\documentclass{article}
\usepackage[T1]{fontenc}
\usepackage{lmodern}
\usepackage{graphicx}
\usepackage{amsmath,amssymb}
\usepackage[a4paper,margin=2cm]{geometry}
\usepackage[absolute,overlay]{textpos}
\setlength{\TPHorizModule}{1mm}
\setlength{\TPVertModule}{1mm}
\usepackage[indonesian]{babel}
\usepackage{hyperref}
\setlength{\emergencystretch}{2em}
% unit: Halaman 38

\begin{document}

% === Halaman 38 (llm) ===

\section*{Keamanan DES}

\begin{itemize}
    \item Keamanan DES ditentukan oleh kunci.
    \item Panjang kunci eksternal DES hanya 64 bit, tetapi yang dipakai hanya 56 bit.
    \item Pada rancangan awal, panjang kunci yang diusulkan IBM adalah 128 bit, tetapi atas permintaan NSA, panjang kunci diperkecil menjadi 56 bit.
    \item Dengan panjang kunci 56 bit akan terdapat $2^{56}$ atau 72.057.594.037.927.936 kemungkinan kunci.
    \item Ruang kunci (\textit{key space}) DES sebanyak $2^{56}$ terlalu kecil, tidak aman karena kunci dapat ditemukan dengan serangan \textit{exhaustive key search} (\textit{brute force attack})
    \item Jika serangan \textit{exhaustive key search} (\textit{brute force attack}) dengan menggunakan prosesor paralel, andaikan dalam satu detik dapat dikerjakan satu juta serangan. Jadi seluruhnya diperlukan 1142 tahun untuk menemukan kunci yang benar.
\end{itemize}

\end{document}
